\chapter*{Conclusion générale et perspectives}
\addcontentsline{toc}{chapter}{Conclusion générale et perspectives}
\markboth{Conclusion générale et perspectives}{Conclusion générale et perspectives}

\myMiniToc{section}{Sommaire}

\vspace{1em}

% =============================================================
\mySectionStar{Problématique étudiée et choix méthodologiques}{}{true}
\label{concl:probl}
% =============================================================

Ce mémoire est parti d'un constat. Aucun modèle actuel de descente d'échelle par apprentissage ne combine la rapidité des réseaux profonds avec une robustesse de principe au changement de distribution, c'est-à-dire au cas où les données rencontrées à l'usage ne ressemblent plus à celles de l'entraînement. Le paradigme à deux étages de \CorrDiff{}~\cite{mardani_corrdiff_2023} stabilise l'entraînement, mais son premier étage de régression générique reste vulnérable à tout déplacement de régime. Notre choix méthodologique a été de \emph{causaliser} ce premier étage en plaçant la causalité dans la topologie du graphe de calcul plutôt que dans une régularisation de la perte. La raison est simple : une régularisation peut être contournée par un réseau suffisamment expressif, alors qu'une dépendance architecturale ne peut, par construction, pas être contournée.

Quatre éléments structurent la contribution, selon les étiquettes annoncées en introduction : (\textbf{C1}) une architecture à deux étages (\emph{two-stage}) où la causalité est inscrite dans la topologie du chemin de calcul ; (\textbf{C2}) un réseau causal récurrent (RCN) intégrant dans une seule cellule le modèle d'équations causales (SCM), la mémoire récurrente (GRU~\cite{cho_gru_2014}) et la contrainte d'acyclicité (DAGMA~\cite{bello_dagma_2022}) présentés aux chapitres~I et~III ; (\textbf{C3}) une interface inter-étages par concaténation de canaux garantissant qu'aucune information ne contourne le chemin causal (propriété O3 vérifiée empiriquement, ratio $0{,}87$) ; (\textbf{C4}) un protocole expérimental à six familles de métriques avec deux diagnostics causaux originaux, à savoir le ratio O3 et le test d'intervention $Q_{\mathrm{int}}$.

% =============================================================
\mySectionStar{Analyse critique des résultats obtenus}{}{true}
\label{concl:analyse}
% =============================================================

Sur la Nouvelle-Zélande (ACCESS-CM2 en distribution ; EC-Earth3 et NorESM2-MM, deux autres modèles climatiques globaux ou GCM, en inter-GCM historique), \Oracle{} réduit le CRPS, l'erreur de prévision probabiliste, de $40\,\%$ en distribution et de $30\,\%$ en inter-GCM par rapport à une \emph{baseline} (modèle de référence) non causale strictement comparable. Le ratio \emph{spread-skill}, qui compare l'incertitude annoncée à l'erreur réellement commise (idéalement 1), passe de $0{,}20$ à $0{,}69$, soit un facteur $3{,}5$ de réduction de la sous-dispersion de l'ensemble, sans atteindre encore la calibration idéale. La distance spectrale RAPSD, qui mesure l'écart de finesse des textures entre cartes produites et cartes observées, chute d'un facteur $450$. Sur deux perturbations contrôlées, \Oracle{} répond avec un signe physiquement cohérent ($Q_{\mathrm{int}}=2/2$), là où la baseline non causale échoue sur l'intervention thermique. Les trois hypothèses de départ reçoivent des verdicts contrastés. Séparer une moyenne déterministe d'un résidu stochastique améliore la stabilité et la calibration : validé sans ambiguïté. La causalité améliore la généralisation d'un modèle climatique à l'autre : validé, mais modérément. Enfin, l'indispensabilité architecturale du graphe (O3) écarte le risque d'un graphe causal purement décoratif : validé par construction.

Plusieurs limites doivent cependant être assumées explicitement :
\begin{itemize}[leftmargin=1.2em, itemsep=2pt]
    \item La matrice $\Adag$ apprise reste à magnitudes uniformes ($Q_{\mathrm{phys}} = 0{,}40$, le score d'accord entre les arêtes apprises et la physique attendue, sur une échelle de 0 à 1) : au vu des données seules, le graphe appris est indiscernable des autres membres de sa classe d'équivalence de Markov (l'ensemble des graphes que des données observationnelles ne permettent pas de départager), conformément à la limite théorique d'identifiabilité~\cite{pearl_causality_2009}. Le DAG appris est une régularisation structurelle inductive, pas une cartographie causale identifiée.
    \item Les magnitudes interventionnelles restent environ trois ordres de grandeur sous la valeur théorique de Clausius-Clapeyron (la loi qui relie réchauffement et humidité disponible) : le signe est correct, la magnitude ne l'est pas.
    \item Le test inter-GCM porte sur le régime historique, distinct d'un vrai test hors-distribution climatique (test sur un scénario d'émissions futur, dit SSP, encore à effectuer).
    \item Aucune donnée africaine n'est testée et la métrique F1@p99 (détection des dépassements du 99\textsuperscript{e} centile) régresse en raison d'une régularisation interne lissante.
    \item Les chiffres reposent sur un seul \emph{seed} (graine aléatoire) d'entraînement (contrainte CPU) et une taille d'ensemble modeste ($K=12$ membres en inter-GCM, $K=64$ en distribution, contre 192~membres pour \CorrDiff{}).
\end{itemize}

% =============================================================
\mySectionStar{Quelques perspectives}{}{true}
\label{concl:persp}
% =============================================================

Cinq axes structurent la suite des travaux, hiérarchisés par horizon temporel :

\textbf{Court terme (1--2 mois).} Une recalibration d'ensemble post-hoc par EMOS~\cite{gneiting_emos_2005} ou BMA (deux méthodes statistiques classiques de correction des ensembles) améliorerait la calibration sans toucher à l'architecture. Nous l'avons écartée ici pour préserver la comparabilité avec la baseline. La validation sur le scénario SSP5-8.5 (le scénario d'émissions le plus intense considéré par le GIEC), exécutable sans modification architecturale sur les données CMIP6 disponibles, transformerait notre test inter-GCM historique en véritable test hors-distribution climatique.

\textbf{Moyen terme (3--6 mois).} Un ancrage prédictif joint de type CASTLE~\cite{kyono_castle_2020} forcerait chaque arête de $\Adag$ à porter un signal effectif, par \emph{fine-tuning} de 25 à 30 époques sans changement d'architecture. En complément, un masque physique sur $\Adag$, via un terme $\lambda_{\mathrm{prior}} \Vert \Adag - \alpha G_{\mathrm{phys}} \Vert^2$, pénaliserait les écarts au couplage quasi-géostrophique attendu (l'équilibre de grande échelle entre vent et pression). Cette piste s'inscrit dans la lignée du \emph{physics-informed ML}~\cite{karniadakis_pinn_2021}.

\textbf{Long terme (1--3 ans).} Le transfert vers l'Afrique subsaharienne reste la motivation finale du travail. Une première étape est une vérification de cohérence par transfert direct, sans ré-entraînement (\emph{zero-shot}), sur les Hautes Terres éthiopiennes ou le Sahel occidental, avec CHIRPS (un produit d'estimation des précipitations combinant satellite et stations au sol, couvrant l'Afrique) comme cible haute résolution. Vient ensuite un ré-entraînement sur des prédicteurs adaptés (mousson, ondes d'Est, zone de convergence intertropicale) et un DAG révisé. Le changement de prior (le graphe causal supposé a priori) est essentiel : le prior néo-zélandais encode un couplage descendant (jet polaire), alors qu'en Afrique équatoriale la convection profonde impose un schéma ascendant (transport intégré d'humidité IVT~$\to$~énergie convective disponible CAPE~$\to$~précipitation de surface SP$_{\mathrm{HR}}$). La réanalyse atmosphérique ERA5~\cite{hersbach_era5_2020} sert de substitut d'initialisation faute de réseau radar dense.

Ce mémoire défend que la causalité utile au downscaling probabiliste n'est pas une propriété décorative émergeant d'une régularisation, mais une contrainte architecturale à placer explicitement dans la topologie du graphe de calcul. Cette approche fait ses preuves en Nouvelle-Zélande aujourd'hui ; elle devra les refaire demain dans les communautés d'Afrique centrale qui ont motivé ce travail.